\section{Sheaves}
\label{section-sheaves}

\noindent
In this section we explain the sheaf condition.

\begin{definition}
\label{definition-sheaf}
Let $X$ be a topological space.
\begin{enumerate}
\item A {\it sheaf $\mathcal{F}$ of sets on $X$} is a presheaf
of sets which satisfies the following additional property: Given
any open covering $U = \bigcup_{i \in I} U_i$ and any collection
of sections $s_i \in \mathcal{F}(U_i)$, $i \in I$ such that
$\forall i, j\in I$
$$
s_i|_{U_i \cap U_j} = s_j|_{U_i \cap U_j}
$$
there exists a unique section $s \in \mathcal{F}(U)$ such that
$s_i = s|_{U_i}$ for all $i \in I$.
\item A {\it morphism of sheaves of sets} is simply a
morphism of presheaves of sets.
\item The category of sheaves of sets on $X$ is denoted
$\Sh(X)$.
\end{enumerate}
\end{definition}

\begin{remark}
\label{remark-confusion}
There is always a bit of confusion as to whether it is
necessary to say something about the set of sections of
a sheaf over the empty set $\emptyset \subset X$.
It is necessary, and we already did if you read the
definition right. Namely, note that the empty set is
covered by the empty open covering, and hence the ``collection
of sections $s_i$'' from the definition above actually form
an element of the empty product which is the final object
of the category the sheaf has values in. In other words,
if you read the definition right you automatically deduce
that $\mathcal{F}(\emptyset) = \textit{a final object}$,
which in the case of a sheaf of sets is a singleton.
If you do not like this argument, then you can just require
that $\mathcal{F}(\emptyset) = \{*\}$.

\medskip\noindent
In particular, this condition will then ensure that if
$U, V \subset X$ are open and {\it disjoint} then
$$
\mathcal{F}(U \cup V) = \mathcal{F}(U) \times \mathcal{F}(V).
$$
(Because the fibre product over a final object is a product.)
\end{remark}

\begin{example}
\label{example-basic-continuous-maps}
Let $X$, $Y$ be topological spaces.
Consider the rule $\mathcal{F}$ which associates to
the open $U \subset X$ the set
$$
\mathcal{F}(U) = \{ f : U \to Y \mid f \text{ is continuous}\}
$$
with the obvious restriction mappings. We claim that
$\mathcal{F}$ is a sheaf. To see this suppose that
$U = \bigcup_{i\in I} U_i$ is an open covering, and
$f_i \in \mathcal{F}(U_i)$, $i\in I$ with
$f_i |_{U_i \cap U_j} = f_j|_{U_i \cap U_j}$ for all $i, j \in I$.
In this case define $f : U \to Y$ by setting $f(u)$
equal to the value of $f_i(u)$ for any $i \in I$ such that
$u \in U_i$. This is well defined by assumption. Moreover,
$f : U \to Y$ is a map such that its restriction to $U_i$
agrees with the continuous map $f_i$. Hence clearly $f$ is
continuous!
\end{example}

\noindent
We can use the result of the example to define constant
sheaves. Namely, suppose that $A$ is a set. Endow $A$ with
the discrete topology. Let $U \subset X$ be an open subset.
Then we have
$$
\{ f : U \to A \mid f\text{ continuous}\}
=
\{ f : U \to A \mid f\text{ locally constant}\}.
$$
Thus the rule which assigns to an open all locally
constant maps into $A$ is a sheaf.

\begin{definition}
\label{definition-constant-sheaf}
Let $X$ be a topological space. Let $A$ be a set.
The {\it constant sheaf with value $A$} denoted $\underline{A}$, or
$\underline{A}_X$ is the sheaf that assigns to an open $U \subset X$
the set of all locally constant maps $U \to A$ with restriction mappings
given by restrictions of functions.
\end{definition}

\begin{example}
\label{example-sheaf-product-pointwise}
Let $X$ be a topological space. Let $(A_x)_{x \in X}$ be
a family of sets $A_x$ indexed by points $x \in X$. We are
going to construct a sheaf of sets $\Pi$ from this data.
For $U \subset X$ open set
$$
\Pi(U) = \prod\nolimits_{x \in U} A_x.
$$
For $V \subset U \subset X$ open define
a restriction mapping by the following rule: An element
$s = (a_x)_{x\in U} \in \Pi(U)$ restricts
to $s|_V = (a_x)_{x \in V}$. It is obvious that this
defines a presheaf of sets. We claim this is a sheaf.
Namely, let $U = \bigcup U_i$ be an open covering.
Suppose that $s_i \in \Pi(U_i)$ are
such that $s_i$ and $s_j$ agree over $U_i \cap U_j$. Write
$s_i = (a_{i, x})_{x\in U_i}$. The compatibility condition implies that
$a_{i, x} = a_{j, x}$ in the set $A_x$ whenever $x \in U_i \cap U_j$.
Hence there exists a unique element $s = (a_x)_{x\in U}$
in $\Pi(U) = \prod_{x\in U} A_x$ with the property that
$a_x = a_{i, x}$ whenever $x \in U_i$ for some $i$. Of course this
element $s$ has the property that $s|_{U_i} = s_i$ for all $i$.
\end{example}

\begin{example}
\label{example-direct-sum-points-not-sheaf}
Let $X$ be a topological space.
Suppose for each $x\in X$ we are given an abelian group $M_x$.
Consider the presheaf $\mathcal{F} : U \mapsto \bigoplus_{x \in U} M_x$
defined in Example \ref{example-direct-sum-points}. This
is not a sheaf in general. For example, if $X$ is an infinite set
with the discrete topology, then the sheaf condition
would imply that $\mathcal{F}(X) = \prod_{x\in X} \mathcal{F}(\{x\})$
but by definition we have $\mathcal{F}(X)
= \bigoplus_{x \in X} M_x = \bigoplus_{x \in X} \mathcal{F}(\{x\})$.
And an infinite direct sum is in general different from an infinite
direct product.

\medskip\noindent
However, if $X$ is a topological space such that every open
of $X$ is quasi-compact, then $\mathcal{F}$ {\it is} a sheaf.
This is left as an exercise to the reader.
\end{example}
